\chapter{Conclusiones}

\section{Valoración personal}

\textbf{Satisfacción con el tutor de la entidad colaboradora.} Satisfacción plena. El Ing. Natanael Pérez Bucio nos dio todo el apoyo a su alcance para la realización de los proyectos planteados, aun cuando uno de ellos (la infraestructura RAG) no llegó a materializarse por las razones ya explicadas. Participó activamente en nuestras reuniones diarias de metodología ágil (el ingeniero está certificado en dicha metodología) y nos apoyó en la evaluación e implementación de las propuestas técnicas del Ing. Rakesh Honnavalli Prabhakara. Fue particularmente valioso que convocara una reunión con los equipos de diseño y de verificación de INNOVA IRV Microelectronics para que mi compañero de prácticas y yo, como expositores del Departamento de Inteligencia Artificial y Ciberseguridad, presentáramos los logros iniciales del proyecto, dándonos así visibilidad dentro de la Fundación. Cabe reafirmar que el Ing. Natanael nos otorgó la libertad creativa necesaria para probar las herramientas y estrategias que nos permitieron desarrollar el proyecto.

\textbf{Satisfacción con el Tutor Interno de la universidad.} Excelente. El profesor Álvaro Calle Cordón mantuvo un contacto permanente conmigo durante la fase previa a la realización de las prácticas, correspondiente a la búsqueda de la entidad colaboradora: me informó de las entidades colaboradoras que iban surgiendo, me dio instrucciones sobre el envío de mi currículum para los distintos procesos de selección, y se comunicó con las entidades para conocer los resultados de selección de los candidatos. Una vez iniciadas mis prácticas, manifestó activamente preocupación por su adecuado desarrollo.

\section{Utilidad como complemento a la formación universitaria}

La utilidad de las prácticas como complemento a mi formación universitaria fue total. Si bien el proyecto que desarrollé no estuvo estrechamente relacionado con el temario cursado en el Máster, supuso un reto que me impulsó a adentrarme en el campo de la Inteligencia Artificial y a aprender, por experiencia directa, cómo abordar un problema en un entorno empresarial real (muy distinto a un problema planteado a nivel académico). Ello implicó enfrentarme a toda la dinámica propia del trabajo en equipo entre distintos departamentos, la intercomunicación entre perfiles técnicos, el manejo efectivo del tiempo, la responsabilidad de cumplir objetivos dentro de plazo, la comunicación profesional en inglés (dado que el asesor técnico externo solo se comunicaba en este idioma, al igual que toda la documentación interna de INNOVA IRV), y la capacidad de gestionar especificaciones de proyecto cambiantes, buscando el consenso entre lo deseado y lo posible en función de los tiempos y las herramientas disponibles.

\section{Utilidad para la futura inserción laboral}

En la misma línea de lo expuesto en el apartado anterior, las prácticas me aportaron una experiencia ya ``probada'' en un entorno profesional real, que resulta directamente transferible a mi futura inserción laboral: la capacidad de trabajar de forma coordinada entre departamentos, de comunicarme eficazmente - incluido en inglés - con perfiles técnicos diversos, de gestionar el tiempo y cumplir objetivos bajo la presión propia de un entorno empresarial, y de adaptarme a especificaciones de proyecto cambiantes buscando soluciones de consenso viables.

\section{Sugerencias de mejora}

\begin{enumerate}
  \item Considero que el departamento de prácticas debería adoptar una mejor planificación para la búsqueda, con la suficiente anticipación, de las mejores entidades colaboradoras posibles.
  \item Sugiero alentar a los estudiantes, desde el inicio del segundo semestre, a buscar y postularse de forma independiente a ofertas de prácticas dentro de la provincia de Málaga, mediante la presentación formal de la asignatura de Prácticas Académicas Externas (PAE), y suministrando toda la información de apoyo que aliente al estudiante a iniciar la búsqueda por cuenta propia.
\end{enumerate}
